% =============================================================================
% 文件: sections/s3_assumptions.tex   第三章 模型假设
% =============================================================================
\section{模型假设}

每条假设都给出“在什么条件下成立”以及“不成立时会带来什么后果”，
以便后续检验章节逐条量化。

\begin{enumerate}[leftmargin=2.4em,itemsep=0.35em]

\item \textbf{长圆柱与一维径向对称假设。}
药材为无限长圆柱，温度与含水率只随径向坐标 $r$ 和时间 $t$ 变化，忽略轴向梯度与端面效应。
\textbf{成立条件}：长径比 $L/(2\Rzero)=6.25$ 较大，且热风沿轴向均匀吹过。
\textbf{后果}：端面额外失水被忽略，等效于低估了约 $8\%$ 的传质面积，
对干燥时长的影响属于系统性偏保守（真实干燥更快）；该误差量级小于传质系数口径
所引入的不确定度。

\item \textbf{局部热力学平衡假设。}
在每一空间点，药材骨架与孔隙水瞬时达到热平衡，可用单一温度 $T(r,t)$ 描述；
水分的迁移以 Fick 型扩散（有效扩散系数 $D$）描述，不单独建立气相压力场。
\textbf{成立条件}：毛细多孔介质在缓慢热风干燥条件下的常用简化，适用于 Biot 数
不是极大的情形（本文 $Bi_h\approx1.25$）。
\textbf{后果}：忽略内部蒸发—冷凝的快速瞬态，可能低估干燥初期表面温度的
“湿球平台”效应。

\item \textbf{物性经验关系的适用范围。}
附录 2 的常数物性、附录 3 与附录 4 的经验公式在 $C\in[0.15,2.55]$、
$T\in[28,51]\ ^\circ\mathrm{C}$ 范围内使用；其中 Arrhenius 因子
$\mathrm{e}^{-3850/T}$ 中的 $T$ 取热力学温度（K）。
\textbf{后果}：公式的声明区间之外（例如含水率低于 $0.15$）的取值属于外推，
不作为结论依据；$\mathrm{e}^{-a/C}$ 在 $C\to0$ 时趋于零，程序中对
$C$ 设置了仅用于避免指数下溢的下限（$10^{-4}$），该下限只在远低于终点判据的
极表层单元被触发。

\item \textbf{含水率变量统一采用干基含水率。}
$C$ 为每千克干物质的含水量（$\mathrm{kg/kg}$）。
\textbf{后果}：附件 1 给出的“烘房水分浓度”被直接用作表面传质边界的驱动势
$C_{\mathrm{air}}$，即认为对流传质系数 $h_m$ 已经把单位换算吸收进去
（见假设 5）。

\item \textbf{对流传质边界条件的驱动势为含水率差。}
表面传质通量写为 $-D\partial C/\partial r=h_m(\Cs-\Cair)$，与附件 2 给出的
$h_m=8\times10^{-7}\ \mathrm{m/s}$ 在量纲上配套。
\textbf{成立条件}：把 $C$ 视为“单位体积干物质中水分质量/干物质质量”的等效浓度，
则 $h_m$ 具有速度量纲。
\textbf{后果}：若实际定义需要显式引入干物质表观密度 $\rho_{\mathrm{d}}$ 或
平衡含水率（吸附等温线），则时间尺度会整体平移；本文在 7.5 节通过
$C_{\mathrm{air}}$ 的整体缩放与 $h_m$ 的反演检验来量化这一不确定度。

\item \textbf{恒温干燥阶段的环境边界为常数。}
$t>14400\ \mathrm{s}$ 后烘房温度与水分浓度保持为常数，基准取附件 1 末段
$3600\ \mathrm{s}$ 的实测均值（$\Tair=49.9957\ ^\circ\mathrm{C}$、
$\Cair=0.049991$）。
\textbf{成立条件}：工艺上恒温干燥阶段通过控制器保持环境稳定。
\textbf{后果}：这是\textbf{模型假设而非数据的唯一推论}；第 7 章用五种口径
（末点、末 $1800/3600/7200\ \mathrm{s}$ 均值、一阶拟合渐近值）量化其影响。

\item \textbf{初始条件均匀。}
$T(r,0)=28\ ^\circ\mathrm{C}$、$C(r,0)=2.55$ 在整个截面上均匀。
\textbf{成立条件}：题面直接给定。

\item \textbf{纯径向收缩，材料点标记法有效。}
问题 4 中长度按题面保持 $25\ \mathrm{cm}$，仅半径按附件 2 变化，
因此采用径向仿射映射 $\xi=r/R(t)$ 标记材料点；这是一种各向异性（纯径向）
收缩假设，而不是三维各向同性收缩。物质坐标下的含水率方程不含附加对流项。
\textbf{后果}：若轴向也发生不可忽略的收缩，或径向变形不是仿射的，需改写为
完整的移动边界/二维模型。

\item \textbf{基准模型不计入汽化潜热。}
热边界只写导热与对流换热，不显式扣除水分汽化所需潜热。
\textbf{理由}：题面未给出可与之配套的汽化潜热与汽化速率耦合关系，
若强行加入会引入与 $h_m$ 口径混淆的第二个未知量。
\textbf{后果}：基准模型可能高估初期升温速率；本文把计入潜热的模型作为
\textbf{扩展模型}单独计算并在 7.5 节报告其影响，不混入基准结论。

\item \textbf{沸腾与辐射换热可忽略。}
药材表面温度始终低于 $60\ ^\circ\mathrm{C}$，无沸腾；烘房壁面辐射折算到
对流换热系数中。

\end{enumerate}
